% Reference: https://en.cppreference.com/w/cpp/container.html
\subsection{stack、queue 与 deque}
\label{stl:queue}
\subsubsection*{原理}
\InlineCode{stack} 后进先出，\InlineCode{queue} 先进先出，均默认用双端队列管理元素。\InlineCode{deque} 可在两端增删，也支持 $O(1)$ 下标访问，但存储不保证连续；中间插入、删除通常为 $O(n)$。
\subsubsection*{使用教程}
栈用 \InlineCode{push/top/pop}，队列用 \InlineCode{push/front/pop}；读取与弹出前先判 \InlineCode{empty()}。\InlineCode{pop()} 不返回弹出元素，先读取再弹出。\InlineCode{deque} 对应 \InlineCode{push_front/push_back}、\InlineCode{pop_front/pop_back}，两端操作为 $O(1)$。
\begin{lstlisting}[style=mtycpp,numbers=none]
stack<int> st;
st.push(3),st.push(5);
if(!st.empty()){int x=st.top();st.pop();}
queue<int> q;
q.push(1),q.push(2);
if(!q.empty()){int u=q.front();q.pop();}
deque<int> d;
d.push_back(4),d.push_front(2);
int x=d[0];
d.pop_back();
q=queue<int>();
\end{lstlisting}
栈、队列和优先队列不提供普通遍历或 \InlineCode{clear()}，可赋为空对象重置。双端队列可维护单调队列；插删后不要依赖先前保存的迭代器。
